% TOP_PROBLEM: {"id": 2652, "title": "Kourovka Notebook Problem 21.143", "category_id": 17, "category": {"id": 17, "name": "group_theory", "display_name": "Group Theory", "description": "Problems about groups, group actions, representations, and related algebraic structures.", "slug": "group-theory", "order_index": 17, "created_at": "2026-07-31T15:26:25.671Z"}, "status": "open", "problem_number": "KOU-21.143", "set_id": 10}
% FIRST_POSTED: 2026-10-01
% ENTRY_KIND: statement_only
% UnsolvedMath ID 2652; source catalogue code KOU-21.143
% !TeX program = lualatex
% Definition and sources only; this file is not an LLM solution attempt.
\documentclass[11pt,a4paper]{article}
\usepackage[margin=25mm]{geometry}
\usepackage{fontspec}
\setmainfont{lmroman10-regular.otf}[BoldFont=lmroman10-bold.otf,
 ItalicFont=lmroman10-italic.otf,BoldItalicFont=lmroman10-bolditalic.otf]
\usepackage{amsmath,amssymb,mathtools}
\usepackage{unicode-math}
\setmathfont{latinmodern-math.otf}
\AtBeginDocument{\let\setminus\smallsetminus}
\usepackage{xurl,hyperref}
\hypersetup{colorlinks=true,urlcolor=blue}
\setcounter{secnumdepth}{0}
\setlength{\parindent}{0pt}
\setlength{\parskip}{5pt}
\setlength{\emergencystretch}{3em}
\begin{document}
\section*{Kourovka Notebook Problem 21.143}\label{2652}
{\small UnsolvedMath ID 2652 · KOU-21.143 · Group Theory · Kourovka Notebook - New Problems, Issue 21}
\subsection{Definitions and mathematical statement}
Thompson's group $F$ is the group, under composition, of orientation-preserving piecewise-linear homeomorphisms of $[0,1]$ with finitely many breakpoints, all breakpoints dyadic rational, and every slope a power of $2$ with integer exponent. A dyadic rational is a number $a/2^b$ with $a\in\mathbb Z$ and $b\ge0$.

For a finitely generated group $G$, take a finite symmetric generating set $S$ and its word metric $d_S$. Let $S^*$ be the set of finite words in $S$, and write $\overline u$ for the group element represented by a word $u$. A language $L\subseteq S^*$ is regular if some finite-state automaton accepts exactly its words.

An automatic structure can be specified by such a regular language $L$ representing every element of $G$, together with a constant $K\ge0$ satisfying the synchronous fellow-traveler condition. For every $u,v\in L$ with $d_S(\overline u,\overline v)\le1$, and every integer $t\ge0$, require
\[
 d_S\bigl(\overline{u_{\le t}},\overline{v_{\le t}}\bigr)\le K,
\]
where $u_{\le t}$ is the prefix of length $\min(t,|u|)$, so a path remains at its endpoint after its word ends. A group is automatic if such $S,L,K$ exist.

Is Thompson's group $F$ automatic in this sense? The requested language need not consist of geodesics or be the usual piecewise-linear normal form. The same finite automaton and fellow-traveler bound must work for all elements and all neighboring represented endpoints.
\subsection{Short English statement}
Does Thompson’s group $F$ admit a regular language of representatives with a uniform synchronous fellow-traveler bound?
\subsection{Sources}
\begin{itemize}
\item[S1] ULAM.AI and the UnsolvedMath Contributors, supplied problems.json, record 2652 (KOU-21.143), Kourovka Notebook - New Problems, Issue 21. Catalogue identity and source transcription; CC BY 4.0.
\url{https://huggingface.co/datasets/ulamai/UnsolvedMath}.
\item[S2] The Kourovka Notebook, 21st edition, version 47 (30 September 2026), new problems, Problem 21.143. Expanded formulation of the supplied catalogue statement.
\url{https://arxiv.org/pdf/1401.0300v47}.
\end{itemize}
\end{document}
